\DocumentMetadata{
  pdfversion=2.0,
  pdfstandard=ua-2,
  lang=en-US,
  tagging=on,
  tagging-setup={math/setup=mathml-SE,tabsorder=structure}
}
\documentclass[11pt,letterpaper]{article}
\usepackage[margin=0.68in,includefoot]{geometry}
\usepackage{newcomputermodern}
\usepackage{amsmath,amscd,mathtools}
\providecommand{\square}{\mdlgwhtsquare}
\usepackage[hidelinks]{hyperref}
\hypersetup{
  pdftitle={Analysis First-Year Exam, January 2015, Part 2},
  pdfauthor={University of Florida Department of Mathematics},
  pdfsubject={Graduate mathematics examination},
  pdfdisplaydoctitle=true,
  bookmarksopen=true
}
\setcounter{secnumdepth}{0}
\setlength{\parindent}{0pt}
\setlength{\parskip}{0.28em}
\setlength{\emergencystretch}{3em}
\setlength{\leftmargini}{2.15em}
\setlength{\leftmarginii}{2.65em}
\setlength{\leftmarginiii}{3em}
\renewcommand{\labelenumi}{\textbf{\theenumi.}}
\pagestyle{empty}
\raggedbottom
\allowdisplaybreaks
\newenvironment{examproblems}[1]{%
  \begin{enumerate}%
  \setcounter{enumi}{#1}%
  \setlength{\topsep}{0.35em}%
  \setlength{\partopsep}{0pt}%
  \setlength{\itemsep}{0.48em}%
  \setlength{\parsep}{0.18em}%
}{\end{enumerate}}
\newenvironment{examsubparts}{%
  \begin{enumerate}%
  \setlength{\topsep}{0.18em}%
  \setlength{\partopsep}{0pt}%
  \setlength{\itemsep}{0.16em}%
  \setlength{\parsep}{0.1em}%
}{\end{enumerate}}
\newenvironment{examinstructions}{%
  \par\begingroup\itshape
}{%
  \par\endgroup\vspace{0.18em}
}
\makeatletter
\renewcommand\section{\@startsection{section}{1}{0pt}%
  {-1.25ex plus -.25ex minus -.1ex}%
  {0.55ex plus .1ex}%
  {\normalfont\large\bfseries}}
\renewcommand{\@maketitle}{%
  \newpage\null\vskip -1.2em%
  {\footnotesize\bfseries
    \href{https://www.ufl.edu/}{University of Florida}, \href{https://math.ufl.edu/}{Department of Mathematics}\par}%
  \vskip 0.18em%
  \tagstructbegin{tag=Title}%
    {\LARGE\bfseries\href{https://gma.math.ufl.edu/exams/analysis-first-year/}{Analysis First-Year Exam}, January 2015, Part 2\par}%
  \tagstructend%
  \vskip 0.35em%
  \tagmcbegin{artifact}%
    \hrule height 1pt%
  \tagmcend%
  \vskip 0.55em%
}
\makeatother
\title{Analysis First-Year Exam, January 2015, Part 2}
\author{}
\date{}
\begin{document}
\maketitle
\thispagestyle{empty}
\begin{examinstructions}
Do exactly 2 problems from Part A and 2 problems from Part B. Work must be presented in a neat and logical fashion in order to receive credit. Do not leave any gaps. When a theorem is used in a proof it must be precisely stated.
\end{examinstructions}
\section{Part A}
\begin{examproblems}{0}
\item
Let \(\sum_{n=0}^{\infty} c_nx^n\) be a power series with finite radius of convergence~\(R\). Assume that the \(c_n\) and~\(x\) are real valued. Let \(\epsilon>0\). State and prove the theorem concerning the uniform convergence of \(\sum_{n=0}^{\infty} c_nx^n\) on \([-R+\epsilon,R-\epsilon]\).
\item
Let \(\{f_n\}\) be a pointwise bounded sequence of complex functions on a countable set~\(E\). Prove that there exists a subsequence \(\{f_{n_k}\}\) of \(\{f_n\}\) such that \(\{f_{n_k}(x)\}\) converges for each \(x\in E\).
\item
Let \(\{f_n\}\) be a uniformly bounded sequence of Riemann integrable functions on \([a,b]\). Let 
\[
F_n(x)=\int_a^b f_n(t)\,dt,\qquad a\leq x\leq b.
\]
 Prove that there exists a subsequence \(\{F_{n_k}\}\) which converges uniformly on \([a,b]\).
\end{examproblems}
\section{Part B}
\begin{examproblems}{0}
\item
Let~\(E\) be a Lebesgue measurable subset of \([0,1]\), with Lebesgue measure \(m(E)\). If \(0<\alpha<m(E)\), show that there exists a Lebesgue measurable subset \(F\subset E\) such that \(m(F)=\alpha\). (Hint: Examine the function \(f(x)=m(E\cap[0,x])\) on \([0,1]\).)
\item
Let \(E\subset\mathbb{R}\) be Lebesgue measurable and suppose \(m(E)<\infty\). It is known that for each \(\epsilon>0\) there is a closed set \(F_\epsilon\subset E\) such that \(m(E\setminus F_\epsilon)<\epsilon\). Show that \(F_\epsilon\) can be chosen to be a compact set.
\item
Let \((X,\Sigma,\mu)\) be a measure space and assume that~\(f\) is~\(\mu\)-integrable. Prove that for \(\epsilon>0\) there exists a \(\delta>0\) such that \(\left|\int_A f\,d\mu\right|<\epsilon\), whenever \(A\in\Sigma\) and \(\mu(A)<\delta\).
\end{examproblems}
\end{document}
